% CS Labs Report --- Laboratory Work No. 1
\documentclass[12pt,a4paper]{article}

\usepackage[utf8]{inputenc}
\usepackage[T1]{fontenc}
\usepackage[english]{babel}
\usepackage{lmodern}
\usepackage[left=2cm,right=1.5cm,top=2cm,bottom=2cm]{geometry}
\usepackage{setspace}
\usepackage{graphicx}
\usepackage{xcolor}
\usepackage{listings}
\usepackage{titlesec}
\usepackage{caption}
\usepackage{float}
\usepackage{amsmath}
\usepackage{amssymb}
\usepackage{fancyhdr}
\usepackage{needspace}
\usepackage[hidelinks]{hyperref}

\graphicspath{{figures/}}
\onehalfspacing
\setlength{\parindent}{1.25cm}
\setlength{\parskip}{0pt}
\setlength{\emergencystretch}{2em}

\definecolor{codebackground}{HTML}{F3F4F6}
\definecolor{codeframe}{HTML}{D1D5DB}
\definecolor{commentgreen}{HTML}{3F6212}
\definecolor{placeholdertext}{HTML}{64748B}

\captionsetup[figure]{name=Fig,labelsep=period,font=small}

\titleformat{\section}
  {\normalfont\bfseries\centering\fontsize{13}{16}\selectfont}
  {\thesection}{0.8em}{\MakeUppercase}
\titleformat{name=\section,numberless}
  {\normalfont\bfseries\centering\fontsize{13}{16}\selectfont}
  {}{0pt}{\MakeUppercase}
\titleformat{\subsection}
  {\normalfont\bfseries\fontsize{12}{15}\selectfont}
  {\thesubsection}{0.7em}{}
\titlespacing*{\section}{0pt}{0pt}{12pt}
\titlespacing*{\subsection}{0pt}{10pt}{6pt}

\pagestyle{fancy}
\fancyhf{}
\fancyhead[L]{\small\textit{Cryptography and Security}}
\fancyhead[R]{\small Laboratory Work No. 1}
\fancyfoot[C]{\thepage}
\renewcommand{\headrulewidth}{0.4pt}
\renewcommand{\footrulewidth}{0pt}
\setlength{\headheight}{15pt}

\lstdefinestyle{go}{
  basicstyle=\ttfamily\small,
  backgroundcolor=\color{codebackground},
  frame=single,
  rulecolor=\color{codeframe},
  breaklines=true,
  columns=fullflexible,
  keepspaces=true,
  showstringspaces=false,
  commentstyle=\color{commentgreen}\itshape,
  morecomment=[l]{//},
  xleftmargin=0.4cm,
  xrightmargin=0.2cm,
  framexleftmargin=0.2cm,
  framesep=5pt,
  aboveskip=8pt,
  belowskip=8pt,
  extendedchars=true,
  literate=
    {Ă}{\u{A}}1 {Â}{\^{A}}1 {Î}{\^{I}}1
    {Ș}{\c{S}}1 {Ț}{\c{T}}1 {Ş}{\c{S}}1 {Ţ}{\c{T}}1
    {ă}{\u{a}}1 {â}{\^{a}}1 {î}{\^{i}}1
    {ș}{\c{s}}1 {ț}{\c{t}}1 {ş}{\c{s}}1 {ţ}{\c{t}}1
    {—}{\textemdash}1
}

\newcommand{\reportfigure}[4][\textwidth]{%
  \begin{figure}[H]
    \centering
    \includegraphics[width=#1]{#2}
    \caption{#3}
    \label{#4}
  \end{figure}
}

\begin{document}

\begin{titlepage}
  \thispagestyle{empty}
  \begin{center}
    \IfFileExists{figures/LOGO_UTM.jpg}{%
      \includegraphics[width=3.4cm]{figures/LOGO_UTM.jpg}\par
      \vspace{0.25cm}
    }{}

    {\bfseries\fontsize{14}{17}\selectfont
      MINISTRY OF EDUCATION, CULTURE AND RESEARCH\par
      OF THE REPUBLIC OF MOLDOVA\par
      Technical University of Moldova\par
      Faculty of Computers, Informatics and Microelectronics\par
      Department of Computer Science and Information Technologies\par
    }

    \vspace{1.4cm}
    {\bfseries\fontsize{16}{19}\selectfont Dinu Plesu, FAF-241\par}

    \vspace{1.4cm}
    {\bfseries\fontsize{36}{40}\selectfont Report\par}
    \vspace{0.25cm}
    {\itshape\fontsize{20}{24}\selectfont Laboratory Work No. 1\par}
    {\bfseries\itshape\fontsize{20}{24}\selectfont Cryptography and Security\par}

    \vspace{1cm}
    {\fontsize{15}{19}\selectfont Topic: Caesar cipher and Caesar cipher with a keyword over the Romanian alphabet\par}
  \end{center}

  \vfill
  \begin{flushright}
    {\fontsize{15}{18}\selectfont Checked by:}\par
    \textbf{Maia Zaica}
  \end{flushright}

  \vfill
  \begin{center}
    {\bfseries\fontsize{14}{17}\selectfont Chi\c{s}in\u{a}u --- 2026}
  \end{center}
\end{titlepage}

\pagenumbering{arabic}
\setcounter{page}{2}
\tableofcontents
\newpage

\section{Task}

\textbf{Task 1.1.} Implement the Caesar cipher for the Romanian alphabet ($n=31$ letters, encoding of Table~2 of the handout) in a programming language of choice. The program must encrypt and decrypt a text given interactively, using the numeric key $k$ ($1$--$30$), and must reject invalid input with a clear message.

\textbf{Task 1.2.} Implement the Caesar cipher with two keys over the same alphabet, keeping all conditions of Task~1.1. Key~1 is the numeric shift; Key~2 is a keyword of at least 7 Romanian letters (including \u{A}, \^{A}, \^{I}, \c{S}, \c{T}) that is converted to uppercase and is used to build a permuted alphabet, which the program must display so the result can be checked by hand.

\section{Theoretical contents}

The Caesar cipher is the simplest monoalphabetic substitution cipher: every letter is replaced by the letter situated $k$ positions later in the alphabet, cyclically. Over an alphabet $A$ of $n$ letters each letter is identified with its code $\mathrm{code}(x)\in\{0,\ldots,n-1\}$; encryption and decryption are

\begin{align*}
  c &= (m+k) \bmod n && \text{(encryption)},\\
  m &= (c-k+n) \bmod n && \text{(decryption)},
\end{align*}
with $m,c\in\{0,\ldots,n-1\}$.

For the Romanian alphabet $n=31$ (Table~2 of the handout), so the key space contains only 30 possible keys; an attacker can recover the plaintext by trying all 30 shifts (brute force) or by comparing the letter frequencies of the ciphertext with the reference frequencies of the language.

In the two-key version (Task~1.2) the second key, a keyword, produces a key-dependent permutation of the alphabet: the keyword letters are written first (duplicates removed) and the remaining alphabet letters follow in their natural order. Encryption again uses a cyclic shift, but now in the permuted alphabet. The key space grows to $30\times(\text{number of admissible keywords})$, yet the cipher remains monoalphabetic and is still broken by frequency analysis.

\section{Implementation}

Both tasks were implemented in Go (files \texttt{task1\_1.go} and \texttt{task1\_2.go}). The alphabet and the letter-to-code map are built once; encryption and decryption differ only in the sign of the shift. The core of Task~1.1 is:

\begin{lstlisting}[style=go,caption={Classic Caesar cipher over the 31-letter Romanian alphabet (\texttt{task1\_1.go})}]
package main

const romanianAlphabet = "AĂÂBCDEFGHIÎJKLMNOPQRSȘTȚUVWXYZ"

func caesarCipher(message []rune, key int, encrypt bool) string {
	alphabet := []rune(romanianAlphabet)
	letterCodes := make(map[rune]int, len(alphabet))
	for code, letter := range alphabet {
		letterCodes[letter] = code
	}

	result := make([]rune, 0, len(message))
	for _, letter := range message {
		code := letterCodes[letter]
		if encrypt {
			code = (code + key) % len(alphabet)
		} else {
			code = (code - key + len(alphabet)) % len(alphabet)
		}
		result = append(result, alphabet[code])
	}
	return string(result)
}
\end{lstlisting}

For Task~1.2 the permuted alphabet is constructed by writing the keyword first and then every Romanian letter not yet used:

\begin{lstlisting}[style=go,caption={Construction of the keyword-permuted alphabet (\texttt{task1\_2.go})}]
func buildPermutedAlphabet(keyword []rune) []rune {
	result := make([]rune, 0, len([]rune(romanianAlphabet)))
	included := make(map[rune]bool)
	for _, letter := range append(keyword, []rune(romanianAlphabet)...) {
		if isRomanianLetter(letter) && !included[letter] {
			result = append(result, letter)
			included[letter] = true
		}
	}
	return result
}
\end{lstlisting}

All inputs are validated interactively: the numeric key must be an integer from 1 to 30, the keyword at least 7 Romanian letters, and the message must contain only Romanian letters and spaces (which are removed). Every violation prints a message stating the allowed values and re-prompts the user.

\section{Results of the work}

\reportfigure{fig1_task1_1.png}{Task 1.1 --- encryption with $k=7$ and the decryption round trip.}{fig:task1-1}

\reportfigure{fig2_validation.png}{Task 1.1 --- invalid key (0, 31) and invalid character \texttt{'2'} are rejected.}{fig:task1-2}

\reportfigure{fig3_task1_2.png}{Task 1.2 --- permuted alphabet built from the keyword \texttt{ROMANIA}, encryption and decryption.}{fig:task1-3}

Together the three runs verify the specification end to end: Fig.~\ref{fig:task1-1} shows that encryption and decryption are exact inverses over the 31-letter alphabet; Fig.~\ref{fig:task1-2} confirms that the input constraints (key in 1--30, only Romanian letters and spaces) are enforced with explicit, re-prompting messages instead of crashing; and Fig.~\ref{fig:task1-3} demonstrates the two-key construction, printing the keyword-derived permuted alphabet ``R O M A N I \u{A} \^{A} B C \ldots'' so that the mapping can be checked by hand against the ciphertext.

\section{Conclusions}

Both ciphers were implemented and validated. The classical Caesar cipher has a key space of only 30 keys, so it offers no real security: a ciphertext is recovered instantly by brute force. Moving the shift into a keyword-permuted alphabet enlarges the key space and hides the simple shift pattern, but the cipher stays monoalphabetic --- each plain letter always maps to the same cipher letter --- so it is still defeated by letter frequency analysis and by the recognisable patterns of frequent words. This is exactly the weakness exploited in Laboratory Work No.~2.

\section*{References}
\addcontentsline{toc}{section}{References}

\begin{enumerate}
  \item Laboratory Work No.~1 --- Caesar Cipher. UTM, FCIM --- handout (\texttt{resources/}).
  \item Table~2. Letter encoding of the Romanian alphabet ($n=31$) --- handout.
  \item W.~Stallings. \textit{Cryptography and Network Security: Principles and Practice}. Classical encryption systems.
\end{enumerate}

\section*{Source Code}
\addcontentsline{toc}{section}{Source Code}

The complete source code of all laboratory works (Go sources for Labs~1--3, Python sources for Labs~4--6, together with the ciphertext and plaintext files used in Lab~2 and Lab~6) is available in the GitHub repository:

\begin{center}
  \url{https://github.com/Prince-of-Nothing/CS-labs.git}
\end{center}

\end{document}
